% =============================================================================
% Section 5 — Conflict Resolution (LWW)
% =============================================================================
\section{Conflict Resolution: Deterministic Last-Write-Wins}
\label{sec:conflict}

When two mutations target the same key and arrive at a node in different orders
on different peers, the nodes must converge to the same winner without
coordination. \pebble{} achieves this through a deterministic total order on
mutations---Last-Write-Wins (LWW).

\subsection{The LWW Total Order}

\begin{definition}[LWW comparison]
\label{def:lww-order}
Define the relation $\lwwlt$ on mutations $a, b \in \MutSet$ as follows:
\[
  a \lwwlt b \;\iff\;
  \begin{cases}
    a.\mts < b.\mts, & \\
    \text{or}\quad a.\mts = b.\mts \;\land\; a.\mnid <_{\mathrm{lex}} b.\mnid, & \\
    \text{or}\quad a.\mts = b.\mts \;\land\; a.\mnid = b.\mnid
               \;\land\; a.\mseq < b.\mseq. &
  \end{cases}
\]
\end{definition}

In prose: the mutation with the \emph{higher timestamp} wins. Ties are broken by
the \emph{lexicographically later node ID}; this is arbitrary but stable across
all nodes because every node applies the same comparison function. A further
tiebreak on \emph{higher sequence} handles the degenerate case of two mutations
from the same node with the same timestamp (possible under fast local writes).

\begin{definition}[Winner]
\label{def:winner}
Given two mutations $a$ and $b$ for the same key, the winner is:
\[
  \win(a, b) =
  \begin{cases}
    b & \text{if } a \lwwlt b, \\
    a & \text{otherwise.}
  \end{cases}
\]
\end{definition}

\subsection{Proof: $\lwwlt$ is a Strict Total Order}

\begin{theorem}[$\lwwlt$ is a strict total order on $\MutSet$]
\label{thm:total-order}
The relation $\lwwlt$ is irreflexive, asymmetric, transitive, and total on
the set of all mutations $\MutSet$.
\end{theorem}

\begin{proof}
We verify each property:

\paragraph{Irreflexivity.}
$a \lwwlt a$ would require $a.\mts < a.\mts$ (false), or $a.\mts = a.\mts$ and
$a.\mnid <_{\mathrm{lex}} a.\mnid$ (false), or all equal and $a.\mseq < a.\mseq$
(false). Hence $a \not\lwwlt a$. \checkmark

\paragraph{Asymmetry (implies antisymmetry).}
Suppose $a \lwwlt b$. We show $b \not\lwwlt a$.
If $a \lwwlt b$ holds because $a.\mts < b.\mts$, then $b.\mts > a.\mts$,
so the first clause of $b \lwwlt a$ is false. The remaining clauses require
$b.\mts = a.\mts$, which is also false. Hence $b \not\lwwlt a$.
The remaining two cases proceed analogously by the asymmetry of $<$ and
$<_{\mathrm{lex}}$. \checkmark

\paragraph{Transitivity.}
Suppose $a \lwwlt b$ and $b \lwwlt c$. We show $a \lwwlt c$ by case analysis
on which clause governs each relation:
\begin{itemize}[leftmargin=1.5em]
  \item If $a.\mts < b.\mts$ and $b.\mts \leq c.\mts$, then $a.\mts < c.\mts$,
        so the first clause of $a \lwwlt c$ holds.
  \item If $a.\mts = b.\mts$ and $b.\mts = c.\mts$ (i.e., all equal), then
        the governing clause is on $\mnid$ or $\mseq$, both of which are
        transitively ordered by the transitivity of $<_{\mathrm{lex}}$ and $<$.
  \item Mixed cases (e.g., $a.\mts < b.\mts$ and $b.\mts = c.\mts$) collapse
        to $a.\mts < c.\mts$ by substitution. \checkmark
\end{itemize}

\paragraph{Totality.}
For any $a \neq b$, either $a.\mts \neq b.\mts$ (handled by the first clause),
or $a.\mnid \neq b.\mnid$ (handled by the second, since $<_{\mathrm{lex}}$ is
total on strings), or $a.\mseq \neq b.\mseq$ (handled by the third, since $<$
is total on integers). The only remaining case is $a = b$, which is excluded.
Hence exactly one of $a \lwwlt b$ or $b \lwwlt a$ holds for $a \neq b$. \checkmark
\end{proof}

\subsection{Convergence Theorem}

\begin{theorem}[Convergence under LWW]
\label{thm:convergence}
Let $n_1, n_2$ be two nodes that have received (in any order) the same set of
mutations $M$ for a key $k$. Both nodes independently compute the same current
value for $k$.
\end{theorem}

\begin{proof}
By \cref{thm:total-order}, $\lwwlt$ is a strict total order on $\MutSet$.
Therefore, the set $\{m \in M : m.\mkey = k\}$ has a unique maximum element
$m^*$ under $\lwwlt$. The function $\win$ selects this maximum deterministically.
Because $\win$ is a function solely of the mutation set---not of arrival
order---both nodes compute the same $m^*$ regardless of the order in which
they processed mutations from $M$. \qed
\end{proof}

\begin{corollary}[Eventual consistency]
\label{cor:eventual}
If the gossip protocol ensures that every mutation created at any node is
eventually delivered to every other node (i.e., the network is eventually
connected), then all nodes eventually agree on the value of every key.
\end{corollary}

\subsection{Implementation}

The comparison function is implemented in \texttt{packages/core/src/store/conflict.ts}:

\begin{lstlisting}[caption={LWW comparison function (TypeScript).}]
export function compare(a: Mutation, b: Mutation): number {
  if (a.timestamp !== b.timestamp) return a.timestamp - b.timestamp;
  if (a.nodeId   !== b.nodeId)   return a.nodeId.localeCompare(b.nodeId);
  return a.sequence - b.sequence;
}

export function winner(incumbent: Mutation, challenger: Mutation): Mutation {
  return compare(challenger, incumbent) > 0 ? challenger : incumbent;
}
\end{lstlisting}

A positive return value indicates the challenger wins; the Store updates
\texttt{currentState[key]} only when the incoming mutation defeats the incumbent.

\subsection{Clock Skew Limitation}

The LWW order as defined relies on wall-clock timestamps. A node whose clock
runs significantly behind can produce mutations that lose to older writes on
other nodes, even if the losing mutation was authored later in real time.
\pebble{} v0.1 accepts this trade-off for simplicity. \Cref{sec:future}
discusses Hybrid Logical Clocks (HLC)~\cite{kulkarni2014hlc} as a path to
causality-aware conflict resolution.
